% =====================================================================
% Preambolo -- Manuale AI and Cybersecurity
% =====================================================================
\usepackage[T1]{fontenc}
\usepackage[utf8]{inputenc}
\usepackage[italian]{babel}
\usepackage{lmodern}
\usepackage[a4paper,margin=2.4cm,headheight=14pt]{geometry}
\usepackage{microtype}
\usepackage{xcolor}
\usepackage{graphicx}
\usepackage{enumitem}
\usepackage{booktabs}
\usepackage{array}
\usepackage{tabularx}
\usepackage{longtable}
\usepackage{multirow}
\usepackage{multicol}
\usepackage{amsmath,amssymb,bm}
\usepackage{listings}
\usepackage[most]{tcolorbox}
\usepackage{tikz}
\usetikzlibrary{mindmap,trees,positioning,arrows.meta,shapes.geometric,shapes.misc,calc,fit,backgrounds}
\usepackage[edges]{forest}
\usepackage{fancyhdr}
\usepackage{hyperref}
\hypersetup{
  colorlinks=true,
  linkcolor=blue!50!black,
  urlcolor=blue!60!black,
  pdftitle={AI and Cybersecurity -- Manuale del corso}
}
\usepackage{bookmark}

\setlist{itemsep=2pt,topsep=4pt}
\renewcommand{\arraystretch}{1.25}
\setcounter{tocdepth}{1}
\setlength{\emergencystretch}{3em}

% ---------------------------------------------------------------------
% Colori e intestazioni
% ---------------------------------------------------------------------
\definecolor{polblu}{HTML}{1F3A5F}
\definecolor{defcol}{HTML}{2E6DA4}
\definecolor{esecol}{HTML}{2E8B57}
\definecolor{codebg}{HTML}{F5F7FA}
\definecolor{codekw}{HTML}{0B5394}
\definecolor{codecm}{HTML}{6A8759}
\definecolor{codestr}{HTML}{A64D79}

\pagestyle{fancy}
\fancyhf{}
\fancyhead[L]{\small\textsc{AI and Cybersecurity}}
\fancyhead[R]{\small\nouppercase{\leftmark}}
\fancyfoot[C]{\thepage}
\renewcommand{\headrulewidth}{0.4pt}

% ---------------------------------------------------------------------
% Riquadri sulle fonti (comuni a tutti i manuali della pipeline)
% ---------------------------------------------------------------------
% Punto che il docente ha indicato come rilevante per l'esame (titolo opzionale)
\newtcolorbox{esame}[1][]{breakable,colback=red!4,colframe=red!60!black,
  fonttitle=\bfseries,title={Per l'esame\ifx&#1&\else: #1\fi}}

% Nota generica (titolo opzionale)
\newtcolorbox{nota}[1][Nota]{breakable,colback=gray!6,colframe=gray!60!black,
  fonttitle=\bfseries,title={#1}}

% Informazione non confermata in aula o ricostruita da noi: da verificare
\newtcolorbox{verifica}{breakable,colback=orange!6,colframe=orange!80!black,
  fonttitle=\bfseries,title=Da verificare}

% Soluzione di quiz, esercizi e simulazioni (titolo opzionale)
\newtcolorbox{soluzione}[1][Soluzione]{breakable,colback=green!4,colframe=green!45!black,
  fonttitle=\bfseries,title={#1}}

% Integrazione dei redattori: contenuto NON detto a lezione né presente nelle slide
\newtcolorbox{integrazione}[1][]{breakable,colback=violet!5,colframe=violet!60!black,
  fonttitle=\bfseries,title={Integrazione -- non dal corso\ifx&#1&\else: #1\fi}}

% Segnaposto per parti ancora da scrivere
\newcommand{\dacompletare}[1]{%
  \begin{tcolorbox}[colback=yellow!8,colframe=yellow!60!black,
    fonttitle=\bfseries,title=Da completare]#1\end{tcolorbox}}

% ---------------------------------------------------------------------
% Riquadri didattici (dal corso)
% ---------------------------------------------------------------------
\tcbset{aicbox/.style={enhanced,breakable,boxrule=0.6pt,arc=2pt,
  left=6pt,right=6pt,top=4pt,bottom=4pt,fonttitle=\bfseries\small,
  before skip=8pt,after skip=8pt}}
\newtcolorbox{definizione}[1][]{aicbox,colback=defcol!5,colframe=defcol,
  title={Definizione\ifx&#1&\else: #1\fi}}
\newtcolorbox{esempio}[1][]{aicbox,colback=esecol!5,colframe=esecol,
  title={Esempio\ifx&#1&\else: #1\fi}}
\newtcolorbox{attenzione}[1][]{aicbox,colback=red!3,colframe=red!45!black,
  title={Attenzione\ifx&#1&\else: #1\fi}}
\newtcolorbox{sintesi}[1][]{aicbox,colback=polblu!4,colframe=polblu,
  title={In sintesi\ifx&#1&\else: #1\fi}}

% ---------------------------------------------------------------------
% Intestazione di una lezione del diario
%   \lezione{data}{titolo}{docenti}{argomenti}
% ---------------------------------------------------------------------
\newcommand{\lezione}[4]{%
  \chapter[#1 -- #2]{#2}%
  \begin{tcolorbox}[colback=blue!3,colframe=blue!50!black,
    fonttitle=\bfseries,title=Scheda lezione]
  \textbf{Data:} #1\\
  \textbf{Docenti:} #3\\
  \textbf{Argomenti:} #4
  \end{tcolorbox}}

% Rimando cliccabile a una lezione del diario: \rif{AAAA-MM-GG}{GG/MM}
\newcommand{\rif}[2]{\hyperref[lez:#1]{lez.~#2}}

% ---------------------------------------------------------------------
% Quiz: domanda numerata + opzioni (a), (b), ...
% ---------------------------------------------------------------------
\newcounter{domanda}
\newenvironment{domanda}[1][]{%
  \refstepcounter{domanda}\par\medskip\noindent
  \textbf{Domanda \thedomanda}\ifx&#1&\else\ \textit{(#1)}\fi\textbf{.}~}%
  {\par}
\newenvironment{opzioni}{\begin{enumerate}[label=(\alph*),nosep,leftmargin=2.2em]}{\end{enumerate}}

% ---------------------------------------------------------------------
% Codice Python
% ---------------------------------------------------------------------
\lstdefinestyle{py}{
  language=Python,
  basicstyle=\ttfamily\footnotesize,
  keywordstyle=\color{codekw}\bfseries,
  commentstyle=\color{codecm}\itshape,
  stringstyle=\color{codestr},
  backgroundcolor=\color{codebg},
  frame=single, rulecolor=\color{black!15},
  numbers=left, numberstyle=\tiny\color{black!45}, numbersep=6pt,
  breaklines=true, breakatwhitespace=true,
  showstringspaces=false, columns=fullflexible, keepspaces=true,
  tabsize=4, xleftmargin=14pt, framexleftmargin=12pt,
  morekeywords={self, True, False, None},
  literate={à}{{\`a}}1 {è}{{\`e}}1 {é}{{\'e}}1 {ì}{{\`i}}1 {ò}{{\`o}}1 {ù}{{\`u}}1
}
\lstset{style=py}

% ---------------------------------------------------------------------
% Macro
% ---------------------------------------------------------------------
\newcommand{\en}[1]{\textit{#1}}             % termine inglese
\newcommand{\code}[1]{\texttt{\small #1}}    % codice inline
\newcommand{\R}{\mathbb{R}}
% Fonte di un capitolo tematico: \fonti{lezioni}{slide}
\newcommand{\fonti}[2]{%
  \noindent{\small\color{black!60}\textbf{Lezioni:} #1 \quad\textbf{Slide:} #2}\par\medskip}

% ---------------------------------------------------------------------
% Mappe concettuali (forest)
% ---------------------------------------------------------------------
\forestset{
  mappa/.style={
    for tree={grow'=0, draw=black!35, rounded corners=2pt, node options={align=left},
      font=\scriptsize, inner sep=2.5pt, l sep=5mm, s sep=1.2mm,
      edge={draw=black!45, thin}, anchor=west, child anchor=west, parent anchor=east},
    forked edges,
    where level=0{fill=polblu!18, font=\small\bfseries, text width=2.4cm}{},
    where level=1{fill=defcol!10, font=\scriptsize\bfseries, text width=3.0cm}{},
    where level=2{fill=white, text width=8.6cm}{},
  },
  mappa3/.style={
    mappa,
    where level=2{fill=esecol!6, text width=3.2cm}{},
    where level=3{fill=white, text width=5.2cm}{},
  }
}
